\section{Prerregulación y protección de celdas de medida (CEN/ECI-41)}
\label{sec:sd4-proteccion-cen41}
Se conservan cuatro LEM IT 200-S para diagnóstico operacional, sin función de seguridad ni permiso de recarga. Se selecciona ajuste de los cuatro TEN 30-2413 a objetivo 16,50 V mediante cuatro resistencias de 1,72 k$\Omega$/0,1\,\%, familia Vishay TNPW0805, entre Trim y $-V_{out}$. La nota OEM vigente identifica para TEN 30-xx13 el par 16,50 V/1,72 k$\Omega$; en el modelo de salida simple los terminales son 4=$+V_{out}$, 5=$-V_{out}$ y 6=Trim. No se usa el valor de otra revisión ni se deduce tolerancia de tensión a partir de tolerancia de la resistencia sin ecuación OEM. La tensión cargada debe aceptarse entre 16,0 y 16,913 V; el extremo superior conserva el cribado propio anterior, no constituye garantía OEM del ajuste. La potencia nominal permanece 30 W: a 16,5 V el límite por potencia equivale a 1,818 A, suficiente para la carga prevista de cada celda, sujeto a derating y arranque reales \parencites[pp.~1--2]{lote41-trim}[pp.~1,3--5]{lote35-ten30}[pp.~1--3]{lote41-tnpw}.

Para aplicar protección positiva independiente a ambas polaridades, se reemplazan expresamente los dos TPS7A3301RGWR negativos de CEN-38 por dos TPS7A4700RGWR adicionales: el total seleccionado es cuatro TPS7A4700RGWR, no seis LDO. Se adoptan cuatro celdas flotantes TEN--LDO--corte, dos por banco. Cada LDO conserva salida local 14,9 V mediante programación ANY-OUT de CEN-38, entrada 10 $\mu$F, salida nominal 47 $\mu$F con capacidad efectiva mínima 20 $\mu$F y NR/SS de 1 $\mu$F. SENSE toma la salida regulada antes del corte, referida al retorno local. La celda positiva tiene retorno local en 0 del banco y salida protegida a +LEM; la celda negativa tiene retorno local a $-\mathrm{LEM}$ y salida protegida a 0 del banco. Sólo estas salidas finales forman la serie bipolar; no se une antes del corte el retorno de la celda negativa a 0, tierra o retorno de campo. Así el controlador recibe tensión positiva respecto de su propio retorno en ambas celdas; la función OEM de protección de polaridad inversa no se interpreta como regulación de una salida negativa. La verificación de retroalimentación, modo común y aislamiento del montaje completo sigue siendo condición de liberación \parencites[pp.~3--6 y 14--18]{lote38-positive}[pp.~2--4 y 7--10]{lote41-ltc}[pp.~5,10--11]{lote35-lem200}.

\textbf{Corte seleccionado.} Cada celda incorpora un Analog Devices LTC4365ITS8\#TRPBF industrial y dos Texas Instruments CSD18540Q5B de 60 V, ocho MOSFET en total. Los N-MOS se conectan en oposición con fuentes comunes, drenaje de entrada al LDO y drenaje de salida a la carga protegida; sus puertas comunes van a GATE. VIN y los divisores UV/OV vigilan la salida LDO anterior al corte; VOUT vigila la salida protegida y GND el retorno local. SHDN queda habilitado por VIN a través de 100 k$\Omega$ en esta etapa de placa. No se presenta esa conexión como secuenciador bipolar. Los terminales FAULT quedan disponibles para futura interfaz aislada, sin conexión directa entre celdas ni a DI24V. No se añade una entrada de PLC aún no seleccionada \parencites[pp.~2--4 y 7--10]{lote41-ltc}[pp.~1--3]{lote41-mos}.

Se selecciona divisor OV de 296 k$\Omega$ arriba (294 k$\Omega$ más 2,00 k$\Omega$ en serie) y 10,0 k$\Omega$ abajo; UV de 253,99 k$\Omega$ arriba (249 k$\Omega$ más 4,99 k$\Omega$) y 10,0 k$\Omega$ abajo. Son valores de familia TNPW0805, 0,1\,\% y 25 ppm/K. Para ambiente de placa 10--50 $^\circ$C se usa error bilateral de resistencia $e=0,001625$, suma conservadora de tolerancia inicial y deriva respecto de 25 $^\circ$C. Con umbral OEM 0,4925--0,5075 V y fuga bilateral 10 nA, los extremos se calculan con $V=V_t[1+R_t(1\pm e)/(R_b(1\mp e))]\pm10^{-8}R_t(1+e)$. Resultan disparo OV 15,0202--15,5814 V y disparo UV 12,9584--13,4420 V. El margen estático mínimo OV frente al absoluto LEM 15,75 V es 0,1686 V; no es margen dinámico demostrado. UV es corte de caída profunda, no protección del límite válido LEM 14,25 V \parencites[p.~4]{lote41-ltc}[pp.~1--3]{lote41-tnpw}[p.~5]{lote35-lem200}.

La histéresis OEM 20--32 mV obliga a evaluar recuperación: para UV resulta 13,4847--14,2894 V y para OV 14,0441--14,9674 V. La exactitud global del TPS7A4700 permite 14,5275--15,2725 V, de modo que algunas combinaciones pueden disparar OV en operación o no rearmar automáticamente. Se exige aceptación conjunta de cada celda cargada sin disparo, incluidos temperatura y transitorios; no se afirma disponibilidad garantizada por nominales ni se estrecha arbitrariamente UV hasta 14,25 V ignorando su histéresis. Queda pendiente el monitor preciso de ventana bipolar y el enclavamiento de parada/arranque conjunto; el estado de medida permanece inválido hasta tener ambas tensiones dentro de 14,25--15,75 V en bornes LEM, estado LEM ON y lecturas frescas coherentes \parencites[pp.~3--4]{lote41-ltc}[p.~6]{lote38-positive}[pp.~5,8,10]{lote35-lem200}.

Se fijan cuatro resistencias de descarga de 1 k$\Omega$/1\,\%/1 W entre salida final y retorno de cada celda, no entre retornos primarios. A 15,75 V consumen como máximo 15,91 mA y 0,2506 W cada una; la corriente por celda prevista sube de 0,56 a 0,57591 A, inferior a 1 A del LDO. Con 16,913 V de entrada y salida 14,5275 V, la pérdida calculada LDO es 1,3738 W, dentro de reserva propia 1,50 W/celda, sujeta a validación térmica instalada. Se reserva resistencia conjunta de MOSFET no mayor que 20 m$\Omega$ a temperatura de placa, como límite propio de aceptación: pérdida 0,00664 W/celda. El máximo OEM 3,3 m$\Omega$ a $V_{GS}=4,5$ V está especificado a 25 $^\circ$C y no se extrapola como máximo térmico. Descarga de una capacidad supuesta de 47 $\mu$F desde 15,75 a 1 V da aproximadamente 0,131 s con 1,01 k$\Omega$; debe medirse la descarga real con todas las capacidades y caminos de retorno, no usarse ese ejemplo como tiempo máximo acreditado \parencites[pp.~1--3]{lote41-mos}[pp.~5--6]{lote38-positive}.

\textbf{Respuesta a fallo y secuencia.} El LTC4365 publica corte rápido máximo 4 $\mu$s sólo con carga de puerta 2,2 nF; cada CSD18540Q5B publica capacidad de entrada máxima 4230 pF, por lo que la pareja no satisface esa condición. Tampoco se transforma el cálculo de carga de puerta en límite de respuesta completo. Un corto IN--OUT del LDO puede elevar su entrada a 16,913 V antes del corte, superior al absoluto LEM: se exige inyección de corto y escalón rápido con primaria aislada, registro osciloscópico de tensión máxima en ambos bornes, corriente, apertura y descarga, y aceptación de pico no superior a 15,75 V. La selección independiente y su umbral estático no prueban aún contención del pico. Pérdida de una celda, corto de MOS, referencia abierta, retorno inadvertido a tierra y alimentación a través de sensores requieren ensayo y coordinación conjunta; el corte de una celda no demuestra desenergización simultánea de ambas. Las conexiones/desconexiones de alimentación LEM se ejecutan con primario separado de partes activas, conforme a su ficha. Estas placas no reciben función SIL ni habilitan recarga \parencites[pp.~3--4]{lote41-ltc}[pp.~1--3]{lote41-mos}[pp.~10--11]{lote35-lem200}.

\textbf{Conciliación energética.} Se conservan cuatro TEN y cuatro PTCB. Los cuatro LDO reemplazan el par positivo/negativo anterior; controladores, MOS, divisores y descargas están aguas abajo de los mismos TEN y caben en su reserva propia de 192 W globales, pendiente de medición a ajuste real. No se agregan otra vez pérdidas LDO, sensores, bus AI o resistencias al balance DC24/AC. Este desarrollo agrega 0 W de reserva upstream. Con ENV-40 incorporado, la base DC es 356,7552 W, y ante pérdida de una QUINT el cálculo de 220,7552 W de la fuente superviviente conserva sólo margen nominal 19,2448 W a condiciones sin derating. No se deduce eficiencia mínima desde el 91\,\% típico TEN. Balance por fase/estado, arranque, refrigeración, bombas, UPS/grupo, autonomía, perfil de recarga y FP/PUE anual permanecen sujetos a los desarrollos y verificaciones correspondientes; RT-06.07, RT-06.08 y RT-06.11 siguen pendientes.
